% Minimal example of the dcptalk style. Build: lualatex example (twice).
\DocumentMetadata{
  lang        = en-GB,
  pdfstandard = ua-2,
  tagging     = on,
}
\documentclass{ltx-talk}
\usepackage[fonts = lm, codedir = .]{dcptalk}
\graphicspath{{images/}}

\title{Talk Title}
\author{David Castro-Perez}
\date{Seminar, 2026-10-19}

\begin{document}

\titleframe{Talk Title}{%
  \authorblock*{David Castro-Perez}{David.CastroPerez@rhul.ac.uk}
    {RHULLogo.png}{Royal Holloway, University of London}\hfill
  \authorblock{A. Coauthor}{coauthor@kent.ac.uk}{KentLogo.png}
    {University of Kent}}{Seminar, 2026-10-19}

\begin{frame}
  \begin{columns}
    \begin{column}{0.4\textwidth}
      \begin{sticky}
        \textbf{The problem}
        \begin{itemize}
          \item one
          \item two
        \end{itemize}
      \end{sticky}
    \end{column}
    \begin{column}{0.55\textwidth}
      \begin{sticky}
        \textbf{Our work}
        \begin{itemize}
          \item three
        \end{itemize}
      \end{sticky}
    \end{column}
  \end{columns}
\end{frame}

\partframe{Part I}{Something}

\begin{frame*}
  \frametitle{Code with a callout}
  \begin{bluebox}
\begin{Verbatim}[commandchars=\\\{\}]
Definition f x := x + 1.\tikzmark{def}
\end{Verbatim}
  \end{bluebox}
  \begin{greenbox}
\begin{Verbatim}
let f x = x + 1
\end{Verbatim}
  \end{greenbox}
  \callout<2->{11,4}{extracts to OCaml}{def}
  \popup<3>(40,60){This box appears on the third slide.}
  \speakernote{Say what \code{f} does; the callout comes on the next
    click. 50\% of the audience will ask about \emph{this}.}
\end{frame*}

\end{document}
